%%%%%%%%%%%%%%%%%%%%%%%%%%%%%%%%%%%%%%%%%%%%%%%%%%%%%%%%%%%%%%%%%%%%%%
% LaTeX Example: Project Report
%
% Source: http://www.howtotex.com
%
% Feel free to distribute this example, but please keep the referral
% to howtotex.com
% Date: March 2011 
% 
%%%%%%%%%%%%%%%%%%%%%%%%%%%%%%%%%%%%%%%%%%%%%%%%%%%%%%%%%%%%%%%%%%%%%%
% How to use writeLaTeX: 
%
% You edit the source code here on the left, and the preview on the
% right shows you the result within a few seconds.
%
% Bookmark this page and share the URL with your co-authors. They can
% edit at the same time!
%
% You can upload figures, bibliographies, custom classes and
% styles using the files menu.
%
% If you're new to LaTeX, the wikibook is a great place to start:
% http://en.wikibooks.org/wiki/LaTeX
%
%%%%%%%%%%%%%%%%%%%%%%%%%%%%%%%%%%%%%%%%%%%%%%%%%%%%%%%%%%%%%%%%%%%%%%
% Edit the title below to update the display in My Documents
%\title{Project Report}
%
%%% Preamble
\documentclass[paper=a4, fontsize=11pt]{scrartcl}
\usepackage[T1]{fontenc}
\usepackage{fourier}

\usepackage[english]{babel}															% English language/hyphenation
\usepackage[protrusion=true,expansion=true]{microtype}	
\usepackage{amsmath,amsfonts,amsthm} % Math packages
\usepackage[pdftex]{graphicx}	
\usepackage{url}
\usepackage{titlesec}
%\usepackage{lstlisting}


%%% Custom headers/footers (fancyhdr package)
\usepackage{fancyhdr}
\pagestyle{fancyplain}
\fancyhead{}											% No page header
\fancyfoot[L]{}											% Empty 
\fancyfoot[C]{}											% Empty
\fancyfoot[R]{\thepage}									% Pagenumbering
\renewcommand{\headrulewidth}{0pt}			% Remove header underlines
\renewcommand{\footrulewidth}{0pt}				% Remove footer underlines
\setlength{\headheight}{13.6pt}
\usepackage[margin=0.5in]{geometry}

%%% Equation and float numbering
\numberwithin{equation}{section}		% Equationnumbering: section.eq#
\numberwithin{figure}{section}			% Figurenumbering: section.fig#
\numberwithin{table}{section}				% Tablenumbering: section.tab#

%% Section Format
\titleformat*{\section}{\LARGE\bfseries}
\titleformat*{\subsection}{\Large\bfseries}
\titleformat*{\subsubsection}{\large\bfseries}
\titleformat*{\paragraph}{\large\bfseries}
\titleformat*{\subparagraph}{\large\bfseries}


%%% Maketitle metadata
\newcommand{\horrule}[1]{\rule{\linewidth}{#1}} 	% Horizontal rule

\title{
		%\vspace{-1in} 	
		\usefont{OT1}{bch}{b}{n}
		\normalfont \normalsize \textsc{University of Purdue\\
      Data Communication and Computer Networks\\
      CS 53600 Fall 2019} \\ [25pt]
		\horrule{0.5pt} \\[0.4cm]
		\huge Assignment 2\\
		\horrule{2pt} \\[0.5cm]
}
\author{
		\normalfont\normalsize Pedro da Costa Abreu Júnior\\ [-5pt]
    \normalsize Student ID 0030878383\\[-5pt]
    \normalsize \today
}
\date{}


%%% Begin document
\begin{document}
\maketitle

\section*{Exercise 3 Discussions}

\begin{enumerate}
\item The experiment B was perfomed running the receiver on pod1-1 and the sender on escher01.

  Running \verb|ping pod1-1.cs.purdue.edu| in escher01 outputs the following:

\begin{verbatim}
escher01 60 $ ping pod1-1.cs.purdue.edu -c 10
PING pod1-1.cs.purdue.edu (128.10.25.201) 56(84) bytes of data.
64 bytes from pod1-1.cs.purdue.edu (128.10.25.201): icmp_seq=1 ttl=60 time=0.872 ms
64 bytes from pod1-1.cs.purdue.edu (128.10.25.201): icmp_seq=2 ttl=60 time=1.08 ms
64 bytes from pod1-1.cs.purdue.edu (128.10.25.201): icmp_seq=3 ttl=60 time=0.962 ms
64 bytes from pod1-1.cs.purdue.edu (128.10.25.201): icmp_seq=4 ttl=60 time=0.964 ms
64 bytes from pod1-1.cs.purdue.edu (128.10.25.201): icmp_seq=5 ttl=60 time=1.08 ms
64 bytes from pod1-1.cs.purdue.edu (128.10.25.201): icmp_seq=6 ttl=60 time=1.10 ms
64 bytes from pod1-1.cs.purdue.edu (128.10.25.201): icmp_seq=7 ttl=60 time=1.04 ms
64 bytes from pod1-1.cs.purdue.edu (128.10.25.201): icmp_seq=8 ttl=60 time=1.05 ms
64 bytes from pod1-1.cs.purdue.edu (128.10.25.201): icmp_seq=9 ttl=60 time=1.24 ms
64 bytes from pod1-1.cs.purdue.edu (128.10.25.201): icmp_seq=10 ttl=60 time=0.995 ms

--- pod1-1.cs.purdue.edu ping statistics ---
10 packets transmitted, 10 received, 0% packet loss, time 9011ms
rtt min/avg/max/mdev = 0.872/1.041/1.245/0.099 ms
\end{verbatim}

Therefore the RTT is 1.117 and K = 2 * 1.041 = 2.082.
Since we can send 1471 bytes in 1.041s in 4s we can send
$1471 * 4 / 1.041 * 1000 \approx 5652257 Bytes$

Using these parameters we get:

\begin{verbatim}
receiver: Total bytes received:	 5652257
receiver: Duplicate Bytes:	 0
receiver: Completion Time:	 4.1446 s
receiver: Speed:		 1363764.8308 bps
\end{verbatim}

Which means that our calculations were right.


The experiment A was performed running the receiver on pod1-1 and the sender on pod2-2.
  Running \verb|ping pod1-1.cs.purdue.edu| in pod2-1 outputs the following:

  \\

\begin{verbatim}
pod2-1 53 $ ping pod1-1.cs.purdue.edu
PING pod1-1.cs.purdue.edu (128.10.25.201) 56(84) bytes of data.
64 bytes from pod1-1.cs.purdue.edu (128.10.25.201): icmp_seq=1 ttl=64 time=0.216 ms
64 bytes from pod1-1.cs.purdue.edu (128.10.25.201): icmp_seq=2 ttl=64 time=0.203 ms
64 bytes from pod1-1.cs.purdue.edu (128.10.25.201): icmp_seq=3 ttl=64 time=0.202 ms
64 bytes from pod1-1.cs.purdue.edu (128.10.25.201): icmp_seq=4 ttl=64 time=0.207 ms
64 bytes from pod1-1.cs.purdue.edu (128.10.25.201): icmp_seq=5 ttl=64 time=0.205 ms
64 bytes from pod1-1.cs.purdue.edu (128.10.25.201): icmp_seq=6 ttl=64 time=0.205 ms
64 bytes from pod1-1.cs.purdue.edu (128.10.25.201): icmp_seq=7 ttl=64 time=0.207 ms
64 bytes from pod1-1.cs.purdue.edu (128.10.25.201): icmp_seq=8 ttl=64 time=0.210 ms
64 bytes from pod1-1.cs.purdue.edu (128.10.25.201): icmp_seq=9 ttl=64 time=0.208 ms
64 bytes from pod1-1.cs.purdue.edu (128.10.25.201): icmp_seq=10 ttl=64 time=0.220 ms
^C
--- pod1-1.cs.purdue.edu ping statistics ---
10 packets transmitted, 10 received, 0% packet loss, time 9222ms
rtt min/avg/max/mdev = 0.202/0.208/0.220/0.012 ms
\end{verbatim}

Therefore the RTT is 0.208 and so K = 2*0.208 = 0.416

Using these parameters, and sending the same file size we get:

\begin{verbatim}
receiver: Total bytes received:	 5652257
receiver: Duplicate Bytes:	 0
receiver: Completion Time:	 0.5988 s
receiver: Speed:		 9439574.6824 bps
\end{verbatim}

As expected, these results shows that running the experiment over the same
network is roughly 6.5 times faster in this case.
Which also meets our expectations, since ping shows that
the connection is 6.5 faster.

\item For the second Part we set dropwhen to 10

  For Experiment B we get the following results:

\begin{verbatim}
receiver: Total bytes received:	 6279947
receiver: Duplicate Bytes:	 628117
receiver: Completion Time:	 4.9360 s
receiver: Speed:		 1272266.3125 bps
\end{verbatim}

As expected, 10\% of the bytes received were duplicates, which means 10% were
dropped.

And we also get a 11\% slowdown, which is roughly what would be expected
since 10\% of the packets were lost.

For experiment A we get:

\begin{verbatim}
receiver: Total bytes received:	 6279947
receiver: Duplicate Bytes:	 628117
receiver: Completion Time:	 1.6722 s
receiver: Speed:		 3755578.1973 bps
\end{verbatim}

\item Now using K of 1.2 RTT we have a timeout of 1.25 for experiment B.
  Running the experiment with the same parameters of item 1 we get

\begin{verbatim}
receiver: Total bytes received:	 5652257
receiver: Duplicate Bytes:	 0
receiver: Completion Time:	 4.2348 s
receiver: Speed:		 1334716.3974 bps
\end{verbatim}

This shows that the numbers did not alter very much, which was expected.
Since the timeout was not triggered to make any difference.

For Experiments A we get very similar results:

\begin{verbatim}
receiver: Total bytes received:	 5652257
receiver: Duplicate Bytes:	 0
receiver: Completion Time:	 0.4973 s
receiver: Speed:		 11365889.80494 bps
\end{verbatim}

\end{enumerate}


%%% End document
\end{document}
